\chapter{Discussion and Limitations}

These findings are directional rather than final, since the analysis does not yet account for intersection delay, time-varying congestion, or all student-specific service requirements (e.g. wheelchair students due to missing data). Nonetheless, they establish a clear proof of concept: better route design can potentially allow \gls{fps} to serve the same students with fewer buses, shorter rides, and more consistent travel times across demographic groups. Additionally, they demonstrate how \gls{fps} can substantially change their bus routing to serve a greater number of students while improving average travel-time outcomes under the modeled assumptions.

To refine these results further, future work will require \gls{fps} to provide a unified student-level dataset linking each student's home location, school, grade, transportation eligibility, current stop assignment, bus assignment, disability status, wheelchair-accessibility requirement where applicable, and any district-specific service constraints. Additionally, I would like to confirm whether routing should use district-approved stops or allow \gls{bra} to generate new stop assignments, and also whether demographic equity constraints should be embedded directly into the \gls{bra} optimization using a strategy similar to \textcite{paviaDesigningEquitableTransit2023}.

As this work is continued, one should define service standards that should govern implementation, including maximum ride time, maximum walk distance to stops, arrival windows, and treatment of students with disabilities and accessibility needs. Future work should replace constant-speed assumptions with calibrated travel times that reflect road type, congestion, and school-period traffic conditions.

\section{Extension to Urban Bus Network Redesigns}

This thesis focuses on a case study for school bus operations in Framingham, Massachusetts, but the underlying co-design technique can be generalized to urban bus networks more broadly. The same monotone compositional reasoning of physical, operational, and policy subsystems applies when planners restructure bus routes, frequencies, and stop patterns across all-day service.

In an urban transit context, several modeling elements change, while the overall structure stays intact. To begin, the demand representation would shift from discrete student assignments to a model of passenger demand, such as \gls{od} matrices, time-of-day boarding patterns, or stop-level ridership patterns. Rather than assigning individual students $m\in\mathcal{M}$ to specific trips, one should embed demand as flows along corridors or between urban activity centers, which can then feed into \gls{mdpi} components that evaluate coverage and travel time distributions for different rider populations.

In addition, stop placement and spacing become explicit design variables, rather than fixed inputs. In this school bus case study, pickup locations are predetermined by student addresses and district policies, and the road network is treated as a base on which to route vehicles. For cities, the co-design problem can include a stop-setting component that trades off walk access distance, stop spacing, and in-vehicle travel time. This would allow one to represent both ``coverage-orientation'' and ``frequency-oriented'' network designs, all with the same formalism.

Furthermore, the temporal structure must expand from single-trip or peak-period service to all-day (and potentially weekend) operations. Instead of enforcing arrival windows around school start and end times, an urban setting would require time-dependent headways, service distribution, and coordination between routes at transfer points, as well as the consideration of multi-modal alternatives. These considerations can be captured by extending the temporal \gls{mdpi} components to include time-of-day frequency decisions and reliability metrics, composing them with fleet and crew scheduling components to model multi-shift operations, and by composing multiple different components together to represent multi-modal operations (possibly by reusing components).

Equity objectives must also shift, from school-specific ride times and eligibility thresholds to accessibility and service quality across diverse neighborhoods and rider groups. The same co-design framework can incorporate equity metrics based on job accessibility, access to healthcare and other key destinations, as well as by implementing constraints on headways, crowding, or travel time disparities between demographic or income groups. In this way, the framework is highly adaptable, providing a unified and formal language for representing a wide variety of transportation planning problems, the main differences residing in how demand, stop placements, temporal patterns, and equity indicators are instantiated and resolved within the shared structure.